\documentclass[11pt]{article}
\usepackage[utf8]{inputenc}
\usepackage{amsmath,amssymb}
\usepackage{hyperref}
\title{Efficient Battery Electrolyte Design: A Resource-Aware Approach}
\author{Anonymous}
\date{}
\begin{document}
\maketitle

\begin{abstract}
This paper studies Battery Electrolyte Design. Grounded in a real, citable literature \cite{ref1,ref2,ref3,ref4,ref5,ref6,ref7,ref8},
we propose a focused method and report \textbf{illustrative} experiments.
All references below are real and verifiable (OpenAlex/arXiv); the reported
numbers are simulated to illustrate intended behaviour.
\end{abstract}

\section{Introduction}
Battery Electrolyte Design is an active research area. This work targets a concrete gap identified from
the recent literature and contributes a method together with an evaluation protocol.

\section{Related Work (real literature)}
The following works, retrieved from public bibliographic APIs, frame our study:
\begin{itemize}
\item Goodenough et al. (2009) — "Challenges for Rechargeable Li Batteries", Chemistry of Materials (cited 10837).
\item Aricò et al. (2005) — "Nanostructured materials for advanced energy conversion and storage devices", Nature Materials (cited 8786).
\item Hwang et al. (2017) — "Sodium-ion batteries: present and future", Chemical Society Reviews (cited 5192).
\item Conway (2013) — "Electrochemical Supercapacitors : Scientific Fundamentals and Technological Applications", CERN Document Server (European Organization for Nuclear Research) (cited 5031).
\item Liu et al. (2019) — "Pathways for practical high-energy long-cycling lithium metal batteries", Nature Energy (cited 3397).
\item Cheng et al. (2012) — "Metal–air batteries: from oxygen reduction electrochemistry to cathode catalysts", Chemical Society Reviews (cited 2661).
\end{itemize}

\section{Method}
We reduce the dominant computational cost via adaptive resource allocation, activating only the components needed per input. We instantiate this idea for Battery Electrolyte Design and analyse its cost/quality trade-off.

\section{Experiments (Illustrative)}
\begin{center}
\begin{tabular}{lcc}
System & Primary Metric & Efficiency \\ \hline
Strong baseline & 0.812 & 1.00$\times$ \\
Ablation & 0.828 & 0.92$\times$ \\
\textbf{Ours} & \textbf{0.851} & \textbf{0.71$\times$}
\end{tabular}
\end{center}
\emph{Metrics simulated to illustrate the intended effect; not measured.}

\section{Conclusion}
We connected a real literature to a concrete method for Battery Electrolyte Design. Future work will
validate the illustrative results with real experiments.

\begin{thebibliography}{9}
\bibitem{ref1} John B. Goodenough, Youngsik Kim. Challenges for Rechargeable Li Batteries. \emph{Chemistry of Materials}, 2009. DOI:10.1021/cm901452z
\bibitem{ref2} A.S. Aricò, Peter G. Bruce, Bruno Scrosati et al.. Nanostructured materials for advanced energy conversion and storage devices. \emph{Nature Materials}, 2005. DOI:10.1038/nmat1368
\bibitem{ref3} Jang‐Yeon Hwang, Seung‐Taek Myung, Yang-Kook Sun. Sodium-ion batteries: present and future. \emph{Chemical Society Reviews}, 2017. DOI:10.1039/c6cs00776g
\bibitem{ref4} Brian E. Conway. Electrochemical Supercapacitors : Scientific Fundamentals and Technological Applications. \emph{CERN Document Server (European Organization for Nuclear Research)}, 2013. 
\bibitem{ref5} Jun Liu, Zhenan Bao, Yi Cui et al.. Pathways for practical high-energy long-cycling lithium metal batteries. \emph{Nature Energy}, 2019. DOI:10.1038/s41560-019-0338-x
\bibitem{ref6} Fangyi Cheng, Jun Chen. Metal–air batteries: from oxygen reduction electrochemistry to cathode catalysts. \emph{Chemical Society Reviews}, 2012. DOI:10.1039/c1cs15228a
\bibitem{ref7} Nian Liu, Zhenda Lu, Jie Zhao et al.. A pomegranate-inspired nanoscale design for large-volume-change lithium battery anodes. \emph{Nature Nanotechnology}, 2014. DOI:10.1038/nnano.2014.6
\bibitem{ref8} Céline Largeot, Cristelle Portet, John Chmiola et al.. Relation between the Ion Size and Pore Size for an Electric Double-Layer Capacitor. \emph{Journal of the American Chemical Society}, 2008. DOI:10.1021/ja7106178
\end{thebibliography}
\end{document}
